\chapter{Discussion, Conclusion and Recommendations}
\label{chap:discussion}

\section{Introduction}
This chapter interprets the results presented in Chapter~\ref{chap:results} in the light of the literature reviewed in Chapter~\ref{chap:literature_review} and the objectives set out in Chapter~\ref{chap:introduction}. Section 5.2 discusses the main findings and what they reveal about the design of the system. Section 5.3 evaluates the extent to which each research objective was achieved, and Section 5.4 states the contributions of the study. Section 5.5 sets out its limitations, Section 5.6 offers recommendations for practitioners, Section 5.7 proposes directions for future work, and Section 5.8 concludes the thesis.

\section{Discussion of Findings}

\subsection{The AI-Score Protects Human Hit Rate When Bots Are Detectable}
The central hypothesis of this study was that feeding a real-time bot-probability score into a cache's admission and eviction logic would preserve the cache's usefulness to human users during heavy scraper activity. The naive-workload results support this hypothesis strongly. DTAC's human hit rate was the same, within measurement error, at 10\% and at 75\% bot traffic (58.29\% and 58.22\%), while LRU lost 11 points over the same range. In effect, detected scrapers were made invisible to the cache: their requests contributed almost nothing to any resource's score, so they never caused an admission (a Cache Pollution Index of exactly zero) and could not displace human content.

The ablation is what allows this effect to be attributed to the AI-score itself rather than to other features of the engine. The no-score variant shares DTAC's ghost cache, decay and eviction tournament and differs only in that it cannot tell bots from humans. The gap between the two grew steadily with the bot share, from 0.33 points at 10\% to 4.33 points at 75\%, while on the stealthy workload, where the router could not distinguish the classes, the gap vanished (between $-0.08$ and $+0.04$ points). These two results together form a natural control: the score helps exactly when, and roughly in proportion to how much, it carries information about which requests are bots.

\subsection{Comparison with Scan-Resistant Algorithms}
\citeA{zhang_rethinking_2026} reported that, under AI scan traffic, ``smarter eviction algorithms can serve as lightweight traffic filters'', with SIEVE, S3-FIFO and ARC degrading far more gracefully than LRU, LFU and FIFO. The present results reproduce that ordering: SIEVE and S3-FIFO outperformed LRU by 12--13 points of human hit rate at 75\% naive bots, and their human hit rate fell by only 2.8 and 4.2 points between 10\% and 75\% bots against LRU's 11. The absolute hit rates are lower than the 70--80\% reported by \citeA{zhang_rethinking_2026} for pure human traffic, because the workloads differ (here the catalog has 100{,}000 objects with a Zipf exponent of 0.8 and the cache holds 10\% of it), but the qualitative behavior is the same. The finding of \citeA{zhang_rethinking_2026} that overall hit ratios fall below 30\% for all policies at a 75\% scan ratio is also visible: at 75\% naive bots no algorithm, DTAC included, exceeded an overall hit rate of 22\%. This is not a failure of the caches; when three quarters of requests are for uniformly scanned objects that will rarely be requested again, most of those requests cannot be hits under any policy. This observation supports the choice of the human hit rate, rather than the overall hit rate, as the primary metric.

Chapter 2 argued that structural scan resistance is ``naive to the origin of the request'' and infers bot-like behavior only from access patterns after the cache is already under pressure. The results quantify the value of adding explicit origin information. On naive bots, DTAC improved on SIEVE by 1.2--3.9 points and on S3-FIFO by 1.0--5.1 points at 10{,}000 entries, and it eliminated cache pollution where SIEVE still allowed 37\% of the cache to be occupied by bot-only content at 75\% bots. It is worth noting, however, that the ablation alone was already roughly on par with SIEVE and S3-FIFO. The engine's two-stage admission, which follows the same intuition as S3-FIFO's small queue and ghost queue \cite{yang_fifo_2023} and TinyLFU's frequency-based admission filter \cite{einziger_tinylfu_2017}, provides most of the structural scan resistance on its own; the AI-score then adds a further, bot-specific layer on top of it.

\subsection{The Engine Behaves Like a Human-Weighted LFU}
Calibration chose a score half-life of four hours, longer than the benchmark logs themselves. Within a replay, therefore, scores barely decay and the engine behaves as a Least Frequently Used policy in which each request is weighted by how human it appears, with recency used only to break ties. Chapter 2 noted that LFU is well suited to static, heavily skewed popularity but suffers from ``frequency pollution'' when popularity changes. The synthetic workloads used here have stationary Zipf popularity, which is precisely the setting that favors LFU. The decay mechanism is what is meant to protect DTAC from frequency pollution on real, drifting workloads, but this study did not evaluate a drifting workload, and the calibrated half-life is likely too long for one. This is an important qualification of the results (Section 5.5).

\subsection{The Limits of Passive Heuristics}
The stealthy workload was designed to show where the heuristics stop working, and the result was unambiguous: with browser headers, residential addresses, a handful of requests per IP and a single favicon fetch, the bots scored marginally \emph{lower} than humans (ROC-AUC 0.485). Each of the four signals is defeated by a simple, cheap change in scraper behavior. The header check tests only for the presence of \texttt{Accept-Language}; the datacenter test is defeated by residential proxies; the velocity and ``Shy Guys'' heuristics both operate per IP and are defeated by rotating addresses, and ``Shy Guys'' additionally needs only one asset fetch per IP. This is consistent with the literature reviewed in Chapter 2: \citeA{ousat_broken_2026} document how modern bots evade even active challenges, and \citeA{hoetzlein_protecting_2025} motivate subnet-level aggregation precisely because per-IP rate limiting is evaded by distributing requests across many addresses. \citeA{van_boxem_shy_2026} present the structural favicon heuristic as a lightweight signal in combination with others, not as a robust one on its own.

The important practical point is that DTAC degraded gracefully when its detector failed. Because an undetected bot is simply treated as a human, DTAC on the stealthy workload behaved exactly like its ablation, which still outperformed LRU and S3-FIFO and tied SIEVE at 75\% bots. A wrong or uninformative score does not make the cache worse than a well-designed score-blind cache; it only removes the additional benefit. The converse risk, penalizing humans, was also small in these experiments: although returning visitors were penalized slightly (mean bot probability 0.030 against 0.021 for new visitors), a human contribution of 0.97 instead of 0.98 per request is immaterial to the cache, and no human crossed a bot probability of 0.5.

\subsection{Sensitivity to Cache Capacity}
DTAC's advantage was not uniform across cache sizes. At 20{,}000 entries it lost to SIEVE at both bot ratios and, at 30\% bots, to S3-FIFO and to its own ablation. Two factors visible in the data explain this. First, the constants were calibrated at 10{,}000 entries, and the calibration check showed that at 20{,}000 a promotion threshold of 0.5 (admitting on the first human request) would have been about 2.3 points better than the calibrated 1.5. With a large cache relative to the catalog, the requirement that an object be requested twice by humans before admission costs more hits than it saves. Second, with 20\% of the catalog cached, the scan-resistant baselines lose comparatively little to bots in the first place, so there is less pollution for the AI-score to prevent. That DTAC's overall hit rate at this capacity fell even below LRU's reflects the same strict admission: objects requested once by a human, or requested mostly by bots, are never admitted, so DTAC forgoes hits that a permissive policy collects. These results show that DTAC's parameters, in particular its promotion threshold, should be tuned to the cache's size relative to its working set, as the literature on adaptive cache management suggests \cite{vysotska_towards_2025}.

\subsection{Cost of the Approach}
The overhead of the approach was modest. The router took 23--39 microseconds at the 99th percentile, far below the 2-millisecond budget set in Chapter~\ref{chap:methodology} and far below the 2.8-millisecond median round trip of even a loopback HTTP request. This confirms the argument of Chapter 2 that composite trust scoring built only from cheap, passive signals can be placed on the critical path of a cache \cite{chandra_composite_2026}. The main cost was memory: DTAC's metadata grew by up to 32.5 MiB against about 5 MiB for LRU, mostly because the ghost cache, bounded at eight times the main capacity, fills up with near-zero-score entries left by bot misses. The additional 28 MiB over LRU amounts to about 11\% of the roughly 256 MiB of payload held in the main cache; it is the price of the long ghost memory that the calibration found beneficial.

\section{Evaluation against the Research Objectives}
Table~\ref{tab:objectives} summarizes the extent to which each specific objective stated in Section 1.3.2 was achieved.

\begin{table}[htbp]
    \centering
    \caption{Status of the research objectives}
    \label{tab:objectives}
    \small
    \begin{tabular}{p{0.3\textwidth}p{0.17\textwidth}p{0.43\textwidth}}
        \toprule
        Objective & Status & Evidence \\
        \midrule
        1. Build a lightweight heuristic detector using passive analysis of header anomalies and structural requests & Achieved against naive bots; not against stealthy bots & Four passive heuristics (velocity, \texttt{Accept-Language}, ``Shy Guys'', datacenter ASN). ROC-AUC 1.000 with no false positives on naive bots; 0.485 on stealthy bots. The ASN signal is correct by construction in the benchmark. \\
        \midrule
        2. Assign a real-time, continuous AI-score in $[0, 1]$ at a bounded per-request cost & Achieved & A continuous log-odds score is computed for every request in $\mathcal{O}(\log K)$ time with bounded memory; $T_{proxy}$ p99 of 23--39 $\mu$s. \\
        \midrule
        3. Develop a DTAC architecture whose admission and eviction are weighted by the AI-score & Achieved & Each request adds $1 - P(\text{bot})$ to its resource's decayed score, which governs both ghost-cache promotion and the eviction tournament. \\
        \midrule
        4. Benchmark DTAC against LRU and scan-resistant algorithms under simulated mixed workloads & Achieved, with a capacity caveat & 350 replays over five seeds with paired confidence intervals. At 10{,}000 entries on naive bots DTAC beat LRU by 6.2--17.1 points, SIEVE by 1.2--3.9 and S3-FIFO by 1.0--5.1; it lost to SIEVE at 20{,}000 entries. \\
        \bottomrule
    \end{tabular}
\end{table}

The primary aim, to design and implement a dynamic heuristic caching architecture that mitigates AI scraper interference and preserves cache efficiency, was therefore achieved for scrapers that exhibit detectable signatures, at cache sizes up to 10\% of the catalog. The expectation stated in Section 3.7.1, that DTAC would maintain a stable human hit rate regardless of the volume of bot traffic, held for detectable bots but not for stealthy ones, against which DTAC's human hit rate declined by a similar amount to SIEVE's and S3-FIFO's.

\section{Contributions of the Study}
The study makes the following contributions:
\begin{enumerate}
    \item \textbf{A working, score-aware cache design.} It presents DTAC, a concrete design in which a continuous bot-probability score weights a resource's popularity score and thereby governs both admission (through a metadata-only ghost cache) and eviction (through a sampled tournament), together with an open, tested Rust implementation of the router, the engine and a reverse proxy built on them.
    \item \textbf{Evidence that origin information adds to structural scan resistance.} Through an ablation that holds the engine fixed, it shows that the AI-score, rather than the admission structure alone, is responsible for most of DTAC's advantage over SIEVE and S3-FIFO, and that this advantage grows with the share of detectable bot traffic.
    \item \textbf{A characterization of when the approach fails.} It shows that the benefit disappears, but performance does not collapse, when bots evade the heuristics, and that the benefit depends on cache capacity and parameter calibration.
    \item \textbf{A reproducible evaluation methodology.} It provides a seeded synthetic workload generator with session-based human traffic, shared human streams across bot ratios, two bot models and stationary mixing, along with a benchmark harness that reports human-specific hit rates, a Cache Pollution Index that does not penalize bot-triggered admissions of content humans use, router classification quality, latency and memory, and that uses held-out calibration and paired statistical comparisons.
\end{enumerate}

\section{Limitations of the Study}
The results should be interpreted in the light of the following limitations.
\begin{enumerate}
    \item \textbf{Synthetic workloads only.} All experiments used synthetic traffic. No labelled real-world traffic was available, so the results show how the system behaves on the modeled traffic, not how it would perform on a real website. In particular, real human browsing is not stationary, and real scrapers are more varied than the two bot models used here.
    \item \textbf{Circularity of the naive bot model.} The naive scrapers were generated with exactly the signatures that the heuristics detect, and the ASN fixture was written to match the generator's datacenter addresses. The perfect detection on this workload therefore demonstrates the mechanism and an upper bound on its benefit, not the accuracy of the detector in practice. For the same reason, the heuristic weights could not be estimated as genuine log-likelihood ratios and were instead calibrated for class separation.
    \item \textbf{Ineffective against stealthy scrapers.} All four heuristics were defeated by bots that imitate browsers and rotate residential IP addresses, which is increasingly how sophisticated scrapers operate \cite{ousat_broken_2026}.
    \item \textbf{Per-IP signals and shared addresses.} The velocity and ``Shy Guys'' heuristics treat every client behind one IP address as a single client. Many users behind a carrier-grade NAT or corporate proxy can resemble a bot, and a bot can hide among them. This effect was not measured.
    \item \textbf{Calibration on one capacity and one generator.} The constants were tuned on a held-out seed, which prevents overfitting to the evaluation seeds but not to the generator itself, and only at a capacity of 10{,}000 entries. DTAC lost to SIEVE at 20{,}000 entries with these constants. The four-hour half-life in particular suits stationary popularity and would likely need to be much shorter for drifting workloads.
    \item \textbf{Entry-based capacity.} Cache capacity was measured in entries, and no evaluated algorithm considers object size. The byte hit rates reported therefore do not reflect a size-aware policy, and results could differ for a cache bounded in bytes.
    \item \textbf{Measurement environment.} The benchmark ran in-process on a single consumer laptop. Hit rates are independent of the hardware, but the latency figures are specific to it and had wide confidence intervals; the networked behavior was validated with only one run, over the loopback interface.
    \item \textbf{Single-node scope.} The system was evaluated as a single cache. Distributed caches and CDN hierarchies, in which per-IP state and popularity would be spread across many nodes, were not considered.
    \item \textbf{Complexity claims.} Contrary to the initial intention of a strictly constant-time router, the bounded per-IP structure updates in $\mathcal{O}(\log K)$ time for a fixed capacity $K$; this cost is independent of traffic volume and was small in practice, but it is not $\mathcal{O}(1)$.
\end{enumerate}

\section{Recommendations}
Based on the findings, the following recommendations are made for practitioners considering score-aware caching:
\begin{enumerate}
    \item \textbf{Start from a scan-resistant engine.} A two-stage admission design (or an algorithm such as SIEVE or S3-FIFO) already recovers most of the hit rate that LRU loses to scrapers, and it is the foundation on which a bot score adds value. Plain LRU should be avoided wherever scraper traffic is significant.
    \item \textbf{Use the bot score as a weight, not a block.} Weighting each request's contribution to popularity by $1 - P(\text{bot})$ gave a large benefit when bots were detectable and degraded gracefully to score-blind behavior when they were not, while the impact of mild false positives on humans was negligible. This makes it a low-risk way to use an imperfect bot detector.
    \item \textbf{Combine many independent signals and expect evasion.} Cheap per-request heuristics catch naive scrapers well but are easily evaded. They should be combined with signals that are harder to imitate, such as subnet-level aggregation \cite{hoetzlein_protecting_2025} or TLS fingerprinting \cite{martinez_llamas_balancing_2025}, and the score should be revisited as scraper behavior changes.
    \item \textbf{Calibrate for the deployment.} The promotion threshold and score half-life should be tuned to the cache's size relative to its working set and to how quickly popularity changes on the site, and the heuristic weights should be reviewed against the site's own human population (for example, reducing the datacenter weight for sites with many VPN users).
    \item \textbf{Budget for metadata.} The ghost cache improves admission decisions but its memory grows with bot traffic up to its bound; operators should size this bound with their memory budget in mind.
\end{enumerate}

\section{Future Work}
Several directions could extend this study:
\begin{enumerate}
    \item \textbf{Evaluation on real traffic.} The most important next step is to evaluate DTAC on real access logs with reliable bot labels, both to measure the true accuracy of the heuristics and to estimate genuine log-likelihood ratios for the weights.
    \item \textbf{Stronger and harder-to-evade signals.} Subnet- and ASN-level aggregation of velocity in the style of Logrip \cite{hoetzlein_protecting_2025}, TLS and HTTP/2 fingerprints, header-ordering consistency, \texttt{robots.txt} interaction \cite{martinez_llamas_balancing_2025}, and lightweight learned classifiers such as BOTracle \cite{kadel_botracle_2025} could be added to the router to address stealthy scrapers.
    \item \textbf{Adaptive parameters.} The promotion threshold and decay half-life could be adapted online, in the spirit of ARC's self-tuning \cite{vysotska_towards_2025}, so that a single configuration performs well across cache sizes and as popularity drifts. Evaluating DTAC under non-stationary workloads is a prerequisite for this.
    \item \textbf{AI-scores on top of other engines.} Since much of DTAC's gain came from the score, the same weighting could be applied to SIEVE or S3-FIFO, for example by withholding the ``visited'' bit or the promotion to the main queue from requests with a high bot probability, to obtain the strengths of both.
    \item \textbf{Size- and cost-aware admission.} Combining the bot score with object size and fetch cost \cite{neglia_access-time-aware_2017} would allow the cache to be bounded in bytes and to optimize byte hit rate and origin load directly.
    \item \textbf{Adversarial and NAT evaluation.} Bot models that adapt to the detector, and human populations behind shared addresses, would show how robust the system is to deliberate evasion and to false positives at scale.
    \item \textbf{Constant-time and distributed state.} A bucket-based stream-summary structure could make the per-IP state update $\mathcal{O}(1)$, and sharing scores and per-IP state across nodes would extend DTAC to multi-node and CDN deployments.
\end{enumerate}

\section{Conclusion}
The rapid growth of AI scraping has broken the assumption of human temporal locality on which traditional web caching relies, and this study set out to restore that locality for human users without partitioning hardware or challenging clients. It designed and implemented in Rust a Dynamic Traffic-Aware Cache that scores every request with a lightweight, passive heuristic router and uses that score to weight each request's contribution to admission and eviction decisions, and it evaluated the system against LRU, SIEVE, S3-FIFO and an ablation of itself over 350 replays of synthetic traffic.

When scrapers were detectable, the approach worked as intended: DTAC kept its human hit rate constant from 10\% to 75\% bot traffic, admitted no bot-only content, and outperformed even modern scan-resistant algorithms by up to about 5 percentage points, at a cost of tens of microseconds per request and a modest amount of metadata memory. When scrapers imitated browsers, the heuristics failed, and DTAC fell back to the performance of a strong score-blind cache rather than below it. Its benefit also depended on cache capacity and on calibration.

The study therefore shows that bot detection and cache management, which have so far been treated separately, can be usefully combined in a single shared cache, and that the quality of the bot score is what determines how much this combination is worth. Making that score robust to evasive scrapers, and validating the approach on real traffic, are the key steps needed to turn this result into a practical defense for web infrastructure.
